\section[Conclusão]{Conclusão}

O projeto meditAMAmente possui a nobre proposta de auxiliar no estudo do impacto da meditação na produção de leite de mães de bebês prematuros, por meio de uma plataforma inovadora e acessível, que promete proporcionar um acompanhamento contínuo sobre os grupos de controle e intervenção participantes do estudo. O projeto visa beneficiar tanto as mães participantes da pesquisa, quanto os pesquisadores ao combinar tecnologia e bem-estar emocional por meio da promoção da prática de meditação. Desde a concepção até a implementação final, cada etapa foi marcada por vitórias e desafios que impulsionaram nosso aprendizado e fortaleceram nossa proatividade e trabalho em equipe.

Nesse projeto, utilizamos tecnologias modernas, arquiteturas robustas e boas práticas de desenvolvimento. Alguns exemplos são: Next.js, Material-UI e MySQL. Isso garantiu que o projeto se tornasse uma aplicação escalável, contando com todas as funcionalidades necessárias, além de vídeos e materiais de apoio para suporte a usuários. A integração de testes automatizados e a organização do código também demonstram o compromisso da nossa equipe com a qualidade e a manutenibilidade do produto desenvolvido, buscando entregar um resultado não somenta satisfatório, mas excelente.

Ao longo de todo o período de desenvolvimento do projeto, foi possível observar um crescimento significativo, tanto técnico e acadêmico, quanto pessoal, de toda a equipe envolvida. A experiência de trabalhar em um contexto real, com prazos e entregas, proporciona uma visão extremamente prática das demandas do mercado, contribuindo para nosso desenvolvimento profissional.

O meditAMAmente alcançou com sucesso todos objetivos planejados e o protótipo entregue atendeu às necessidades dos stakeholders e, consequentemente, do público-alvo. Acredito que isso se dá, principalmente, pelo potencial transformador, tanto na esfera acadêmica quanto social, que o projeto demonstrou, despertando o interesse de todos os envolvidos e, consequentemente, contribuindo para o sucesso do trabalho.

Em suma, o meditAMAmente, mais do que um trabalho técnico da faculdade, foi uma experiência extremamente enriquecedora, da qual fiquei muito feliz e orgulhosa de ter participado. O empenho e os constantes esforços de toda a equipe resultaram em um produto relevante para a sociedade. Acredito que eu também pude contribuir significativamente para o projeto, de modo que fiquei bastante satisfeita com meu desempenho. Durante esse semestre, pude desenvolver minhas habilidades em diversas áreas diferentes como banco de dados, desenvolvimento front-end, integração de ambientes, arquitetura de software, coolaboração, trabalho em equipe, tutoria, entre muitas outras. Logo, os resultados atingidos por meio da atuação nesse período na \ac{ages} que representam o compromisso, dedicação, organização e companheirismo desse time maravilhoso, no qual cada integrante deu o seu máximo para fazer uma diferença positiva no mundo por meio da tecnologia.
